\documentclass[10pt,a4paper]{article}
\usepackage[T1]{fontenc}
\usepackage{lmodern,microtype}
\usepackage[margin=25mm]{geometry}
\usepackage{amsmath,amssymb,amsthm}
\usepackage[colorlinks=true,linkcolor=blue,urlcolor=blue]{hyperref}
\hypersetup{pdftitle={Systematic and random components of propagated error}}
\newcommand{\R}{\mathbb{R}}
\newcommand{\E}{\mathbb{E}}
\newcommand{\norm}[1]{\left\lVert#1\right\rVert}
\DeclareMathOperator{\tr}{tr}
\DeclareMathOperator{\Cov}{Cov}
\newtheorem{proposition}{Proposition}
\title{Systematic and random components of propagated error}
\author{}
\date{15 November 2026}

\begin{document}
\maketitle
\noindent\textbf{Status:} established moment identities with a conditional
first-order network interpretation. No novelty or empirical claim is made.

\section{Context and notation}
We study how local numerical errors combine at a common receiving output.
A local error is the difference between an approximate value and its reference
value, measured at the same point in the computation. It may arise from
quantization, rounding, or physical noise; no particular mechanism is assumed.

There are $L\geq1$ error locations, indexed by $o,p\in\{1,\ldots,L\}$.
At location $o$, the error is a column vector $\delta_o\in\R^{n_o}$;
$n_o\geq1$ is its number of scalar entries. Tensors are flattened using a fixed
ordering. The common receiving output has $m\geq1$ entries.
Here $\R$ denotes the real numbers, $T$ denotes transpose, and $I_n$ is the
$n\times n$ identity matrix. All transport arithmetic below is over real
numbers; any additional arithmetic error must be included as another error
contribution if it is to be covered.

The symbol $\E$ denotes expectation over one explicitly specified ensemble
of errors. Every mean and covariance below uses that same ensemble and is
assumed finite. For random column vectors $a,b$, define
$\Cov(a,b)=\E[(a-\E a)(b-\E b)^T]$ and $\Cov(a)=\Cov(a,a)$.
For a scalar random variable $z$, $\operatorname{Var}(z)=\E[(z-\E z)^2]$.
The root-mean-square (RMS) vector magnitude is $\sqrt{\E\norm{a}_2^2}$;
it is not divided by the vector length. For a vector $a$,
$\norm{a}_2=\sqrt{a^Ta}$ is its Euclidean norm. For a matrix $A$,
$\tr(A)$ is its diagonal sum when $A$ is square,
$\norm{A}_F=\sqrt{\tr(AA^T)}$ is its Frobenius norm, and
$\norm{A}_2=\sup_{\norm{u}_2=1}\norm{Au}_2$ is its spectral norm.

\section{Two components and the transport being studied}
For a declared ensemble of local errors at operator $o$, write
\begin{equation}
 \delta_o=\mu_o+\xi_o,\qquad
 \mu_o=\E\delta_o,\qquad \E\xi_o=0,\qquad
 C_{op}=\E[\xi_o\xi_p^T],\quad\Sigma_o=C_{oo}.
 \label{sr:decomposition}
\end{equation}
The systematic component $\mu_o\in\R^{n_o}$ is the mean error;
$\xi_o=\delta_o-\mu_o\in\R^{n_o}$ is the centered fluctuation.
The matrix $C_{op}\in\R^{n_o\times n_p}$ is the cross-covariance between
locations $o$ and $p$; $\Sigma_o\in\R^{n_o\times n_o}$ is the covariance
at location $o$. This decomposition is exact whenever second moments exist.
The expectation must identify its ensemble: repeated physical noise, rounding
randomness, or a specified input population. Deterministic quantization of a
fixed input does not automatically produce random, unbiased noise.

Let $M_o\in\R^{m\times n_o}$ be a fixed matrix mapping local error $\delta_o$
to the common output space. Define the total transported error $d\in\R^m$ by
\begin{equation}
 d=\sum_{o=1}^L M_o\delta_o.
 \label{sr:transport}
\end{equation}
For a network example, let $h_0\in\R^{n_0}$ be a fixed reference input and
$h_o=f_o(h_{o-1})\in\R^{n_o}$ its reference intermediate values, with
$n_L=m$. The differentiable map $f_o:\R^{n_{o-1}}\to\R^{n_o}$ is operator
$o$. Consider the perturbed computation
$\widetilde h_0=h_0$ and
$\widetilde h_o=f_o(\widetilde h_{o-1})+\delta_o$:
the local error is injected at the \emph{output} of that operator, and may
depend on preceding errors. Its reference Jacobian is
$J_o=\left.\partial f_o/\partial h\right|_{h=h_{o-1}}
\in\R^{n_o\times n_{o-1}}$, where $h$ is the input argument of $f_o$.
The first-order transport is $M_o=J_L\cdots J_{o+1}$ for $o<L$, and
$M_L=I_m$. The actual final-output error is
$\Delta=\widetilde h_L-h_L\in\R^m$; write $\Delta=d+R$, with
$R=\Delta-d\in\R^m$ the linearization remainder containing nonlinear effects.
The identities below are exact for $d$, not a universal exact law for $\Delta$.
If inputs vary, the Jacobians also vary: condition on the reference input or
retain their joint dependence when averaging.
This is a chain-rule specialization of first-order uncertainty propagation;
see~\cite[Sections 5.1--5.2, Eqs.~(10), (13)]{jcgm2008gum}.

\section{Mean direction and covariance propagate differently}
\begin{proposition}[General second-moment identity]
With fixed $M_o$ and finite second moments,
\begin{align}
 \E\norm{d}_2^2
 &=\left\lVert\sum_o M_o\mu_o\right\rVert_2^2
   +\sum_o\tr(M_o\Sigma_o M_o^T) \notag\\
 &\quad+2\sum_{o<p}\tr(M_o C_{op}M_p^T).
 \label{sr:general}
\end{align}
Here $\sum_o$ means $\sum_{o=1}^L$, and $\sum_{o<p}$ sums once over every
pair $1\leq o<p\leq L$.
\end{proposition}
\begin{proof}
Write $d=\bar d+\eta$, where $\bar d=\E d=\sum_oM_o\mu_o\in\R^m$ is its
mean and $\eta=d-\bar d=\sum_oM_o\xi_o\in\R^m$ its centered fluctuation.
Since $\E\eta=0$,
$\E\norm{d}_2^2=\norm{\bar d}_2^2+\tr\Cov(\eta)$.
Expanding $\Cov(\eta)$ gives~\eqref{sr:general}.
\end{proof}
The covariance and quadratic-moment identities used in this proof are standard;
see~\cite[Section 6.2, Eqs.~(313)--(314), (327)--(329)]{petersen2012matrixcookbook}.

The systematic contributions add as vectors before taking their norm.
Their transported directions may reinforce, oppose, or be annihilated.
Fluctuations propagate through covariance, including correlations between
operators. Zero mean alone does not remove the cross-covariance terms.

If fluctuations at different operators are uncorrelated ($C_{op}=0$ for
$o\ne p$), their contributions add in squared norm. Independence is sufficient
but stronger than necessary. With the additional isotropic assumption
$\Sigma_o=\sigma_o^2 I_{n_o}$, where $\sigma_o\geq0$ is the common standard
deviation of coordinates at location $o$, the expression reduces to
\begin{equation}
 \E\norm{d}_2^2=
 \left\lVert\sum_o M_o\mu_o\right\rVert_2^2
 +\sum_o\sigma_o^2\norm{M_o}_F^2.
 \label{sr:isotropic}
\end{equation}
This is the systematic/centered-fluctuation split under explicit assumptions.
Isotropy and zero cross-covariance each require justification; a shared scale
or reused hardware can introduce correlations between contributions.

\section{The two gains require consistent normalization}
For one operator, suppress its index: $n$ is its local error dimension,
$M\in\R^{m\times n}$ its transport, $\mu\in\R^n$ its mean error,
$\xi\in\R^n$ its centered fluctuation, and $\Sigma=\Cov(\xi)$ its covariance.
For $\mu\ne0$, define its magnitude $\epsilon_s=\norm{\mu}_2$ and unit
direction $u=\mu/\epsilon_s$. Its directional gain $g_s(u)$ is defined by
\begin{equation}
 g_s(u)=\norm{Mu}_2,\qquad \norm{M\mu}_2=\epsilon_s g_s(u).
 \label{sr:directional-gain}
\end{equation}
If $\mu=0$, its transported mean is zero and no mean direction is needed.
For isotropic fluctuations, $\sigma$ is a \emph{per-coordinate standard
deviation}, meaning $\Sigma=\sigma^2 I_n$ with $\sigma\geq0$, whereas
$r=\sqrt{\E\norm{\xi}_2^2}=\sigma\sqrt n$ is the RMS
\emph{vector norm}. Consequently,
\begin{equation}
 \sqrt{\E\norm{M\xi}_2^2}
 =\sigma\norm{M}_F
 =r\,\frac{\norm{M}_F}{\sqrt n}.
 \label{sr:isotropic-gain}
\end{equation}
Thus the gain relative to RMS vector error is $\norm{M}_F/\sqrt n$.
Comparing $\norm{Mu}$ with unnormalized $\norm{M}_F$ at supposedly equal
vector-error magnitude mixes conventions. The spectral norm is a worst-direction
bound; $\norm{M}_F/\sqrt n\leq\norm{M}_2$.
For general covariance with $\tr\Sigma>0$, define the RMS gain
$g_{\Sigma}$ as output fluctuation RMS divided by input fluctuation RMS:
\begin{equation}
 g_{\Sigma}=\sqrt{\frac{\tr(M\Sigma M^T)}{\tr\Sigma}}.
 \label{sr:covariance-gain}
\end{equation}

\paragraph{Equal-size example.}
Set $n=2$ and $m=1$. Take a deterministic error $\mu=a(1,1)^T$, where $a>0$
is an arbitrary scalar amplitude, and compare it with a separate zero-mean
random error $\xi$ having covariance $\Sigma=a^2I_2$.
Both input errors have RMS vector magnitude
$\sqrt2a$. With $M_+=\begin{bmatrix}1&1\end{bmatrix}$, the systematic output
magnitude is $2a$ and the fluctuation RMS is $\sqrt2a$. With
$M_-=\begin{bmatrix}1&-1\end{bmatrix}$, the systematic output vanishes while
the fluctuation RMS remains $\sqrt2a$. A bias is not intrinsically more harmful:
its direction and the receiving map matter.

\section{When linear and square-root growth hold}
Suppose $M_o\mu_o=v$ for every $o$, where $v\in\R^m$ is a fixed vector.
Suppose each transported centered fluctuation has the same RMS norm
$s=\sqrt{\E\norm{M_o\xi_o}_2^2}>0$, with zero cross-covariance between
different transported fluctuations. Then
\begin{equation}
 \left\lVert\sum_oM_o\mu_o\right\rVert_2=L\norm{v}_2,
 \qquad
 \sqrt{\E\left\lVert\sum_oM_o\xi_o\right\rVert_2^2}=\sqrt L\,s.
 \label{sr:depth}
\end{equation}
Only in this equal-contribution, aligned-mean case does mean error dominate
when $\norm{v}_2/s>1/\sqrt L$. These are conditional scaling laws, not a
universal bias-versus-noise ordering.

Two counterexamples delimit the claim. Transported means $v$ and $-v$ cancel.
Conversely, if every transported fluctuation is the same random vector $Zv$,
where $Z$ is one shared scalar random variable with $\E Z=0$ and $\E Z^2=1$,
their sum has RMS norm $L\norm{v}_2$.
It is zero-mean but perfectly correlated and therefore has linear RMS growth.
The word \emph{coherent} describes alignment/correlation after transport;
nonzero local mean does not itself guarantee such alignment across layers.

The same accounting applies to a finite dot product. Let $K\geq1$ be the
number of reduction terms, $w=(w_1,\ldots,w_K)^T\in\R^K$ the fixed weights,
and $e=(e_1,\ldots,e_K)^T\in\R^K$ the random errors in their corresponding
operands. Define the scalar output error $z=\sum_{k=1}^K w_ke_k=w^Te$.
Suppose $\E e=\beta\mathbf1_K$ and $\Cov(e)=\sigma^2I_K$, where
$\beta\in\R$ is the common coordinate mean, $\mathbf1_K$ is the vector of
$K$ ones, and $\sigma\geq0$ is the common coordinate standard deviation.
Then
\begin{equation}
 \E z=\beta\sum_{k=1}^K w_k,\qquad
 \sqrt{\operatorname{Var}(z)}=\sigma\sqrt{\sum_{k=1}^K w_k^2}.
 \label{sr:reduction}
\end{equation}
The familiar $K\beta$ versus $\sqrt K\,\sigma$ follows for $w_k=1$;
arbitrary weights can change the comparison or cancel its mean term.
These additive-error RMS identities differ from the probabilistic
backward-error bounds for computed inner products
in~\cite[Section 3.1]{higham2019probabilistic}.

\section{What the distinction supports}
Mean, variance, and relevant correlations describe different parts of an error;
one mean-squared-error (MSE) number generally cannot identify their downstream
effect. Here MSE means a squared error averaged over the declared ensemble,
with any coordinate normalization stated explicitly. A practical
estimator must specify its ensemble, normalization, and available observations.
These identities do not supply downstream Jacobians for free or establish an
online control algorithm. Final task quality requires separate evaluation of
the actual perturbed computation, including nonlinear and input-dependent effects.

The companion note
\href{activation-error-direction.pdf}{\emph{Activation-error direction and
matrix-output error}} gives the deterministic, single-matrix counterpart.
Both notes are mathematical background, with no claim that cancellation is
frequent in trained networks or that a particular compensation policy helps.

\section{Literature and provenance}
\paragraph{Established identities.}
Petersen and Pedersen~\cite{petersen2012matrixcookbook} collect the moment
identities used in~\eqref{sr:general}; Section 10.4 supplies the matrix norms.
Equations~\eqref{sr:isotropic}--\eqref{sr:covariance-gain},
\eqref{sr:depth}, and~\eqref{sr:reduction} are algebraic specializations
derived here. The examples and counterexamples are constructed illustrations.

\paragraph{First-order interpretation.}
JCGM~\cite{jcgm2008gum} gives the correlated-input uncertainty framework.
Its Annex E.3.2 distinguishes uncertainty from realized error; applying
the framework to local numerical errors requires the declared ensemble
and linearization, rather than a neural-network accuracy guarantee.

\paragraph{Assumptions and rounding models.}
Higham and Mary~\cite[Model 2.1; Sections 3.1, 4.2]{higham2019probabilistic}
analyze bounded, independent, mean-zero rounding errors and give failure cases.
They do not establish isotropic neural quantization noise or the depth law
in~\eqref{sr:depth}. Isotropy, uncorrelated errors, and aligned equal
contributions are chosen assumptions here, requiring separate validation.

\renewcommand{\refname}{References}
\begin{thebibliography}{9}
\bibitem{petersen2012matrixcookbook}
K.~B. Petersen and M.~S. Pedersen.
\emph{The Matrix Cookbook}. Version 15 November 2012.
Sections 6.2 and 10.4.
\href{https://www2.imm.dtu.dk/pubdb/edoc/imm3274.pdf}{Author PDF}.

\bibitem{jcgm2008gum}
Joint Committee for Guides in Metrology.
\emph{Evaluation of measurement data---Guide to the expression of
uncertainty in measurement}. JCGM 100:2008, corrected version 2010.
Sections 5.1--5.2 and Annex E.3.2.
\href{https://doi.org/10.59161/JCGM100-2008E}{doi:10.59161/JCGM100-2008E}.

\bibitem{higham2019probabilistic}
N.~J. Higham and T.~Mary.
A new approach to probabilistic rounding error analysis.
\emph{SIAM Journal on Scientific Computing}, 41(5):A2815--A2835, 2019.
\href{https://doi.org/10.1137/18M1226312}{doi:10.1137/18M1226312}.
\href{https://eprints.maths.manchester.ac.uk/2731/1/paper.pdf}{Author manuscript},
18 March 2019; Model 2.1 and Sections 3.1, 4.2.
\end{thebibliography}
\end{document}
